\section{Diffusione dell'attacco}\label{sec:Diffusione}

\subsection{Diffusione iniziale}\label{subsec:diffusioneIniziale}

ILOVEYOU era stato inizialmente progettato per funzionare solo a Manila; questa limitazione fu poi rimossa, permettendo al worm di diffondersi in tutto il mondo \cite{white2020}. Le prime segnalazioni ai produttori di antivirus risalgono alla mattina del 4 maggio 2000 \cite{whalen2000}. Seguendo i fusi orari, il worm si spostò verso ovest man mano che i dipendenti iniziavano la giornata e aprivano la posta: prima Hong Kong nel pomeriggio, poi l'Europa, dove era mattina, e infine gli Stati Uniti al risveglio \cite{gao2000iloveyou,bbc2000loveyou,guardian2000lovebug}.

\subsection{Diffusione in enti governativi e aziende}

Il worm paralizzò la posta elettronica di molte organizzazioni, che in diversi casi dovettero spegnere i server per il volume di messaggi. Secondo il rapporto del GAO \cite{gao2000iloveyou}, furono colpiti, tra gli altri:

\begin{itemize}
    \item in Belgio, il sistema bancario nazionale;
    \item nel Regno Unito, i server di posta della \textit{House of Commons}, spenti per due ore;
    \item negli Stati Uniti, la maggior parte delle agenzie federali, tra cui il \textit{Dipartimento di Giustizia}, il \textit{Dipartimento del Lavoro} e la \textit{Social Security Administration}, oltre a \textit{CIA} e \textit{Dipartimento della Difesa}; alcuni file della \textit{NASA} non furono recuperabili nemmeno dai backup.
\end{itemize}

La sola \textit{Veterans Health Administration} ricevette circa 7 milioni di email del worm, e servirono 240 ore di lavoro per risolvere i problemi \cite{gao2000iloveyou}.

\subsection{Dispositivi colpiti e danni}

Le prime stime, riportate da Whalen \cite{whalen2000}, parlavano di almeno 650 siti e oltre 500.000 sistemi coinvolti. Dopo circa dieci giorni si stimò che il worm avesse infettato circa \textbf{50 milioni} di computer, circa il \textbf{10\%} di quelli connessi a Internet \cite{winder2020}.

Il bersaglio erano i sistemi Windows: Mac e UNIX non potevano eseguire lo script, ma potevano comunque inoltrare, manualmente o tramite inoltri automatici, il file infetto ad altri utenti \cite{whalen2000}.

\subsection{Rimozione}

Non bastava cancellare il file ricevuto: il worm si era copiato nel sistema e registrato all'avvio (Sezione~\ref{subsec:fase1}), quindi serviva uno strumento dedicato \cite{whalen2000}. Strumenti di rimozione e aggiornamenti delle firme furono rilasciati, tra gli altri, da \textbf{McAfee}, \textbf{F-Secure}, \textbf{Symantec}, \textbf{Sophos} e \textbf{Trend Micro}, che dovettero distribuirli molto più rapidamente del solito \cite{ebsco2023}.

La parte più difficile fu il recupero dei file sovrascritti (Sezione~\ref{subsec:fase3}) da backup, quando disponibili: la maggior parte del costo complessivo fu dovuta ai giorni di lavoro dei tecnici, più che al software di rimozione.
